\documentclass[11pt,a4paper]{article}
\usepackage[T1]{fontenc}
\usepackage{lmodern,amsmath,amssymb,amsthm,mathtools,booktabs,microtype}
\usepackage[margin=27mm]{geometry}
\usepackage[colorlinks=true,linkcolor=blue,citecolor=blue,urlcolor=blue]{hyperref}
\setlength{\emergencystretch}{2em}
\newtheorem{theorem}{Theorem}[section]
\newtheorem{lemma}[theorem]{Lemma}
\newtheorem{proposition}[theorem]{Proposition}
\newtheorem{corollary}[theorem]{Corollary}
\newcommand{\Q}{\mathbb Q}\newcommand{\Z}{\mathbb Z}
\DeclareMathOperator{\Tr}{Tr}\DeclareMathOperator{\Nm}{N}
\title{Quadratic reweightings of native elliptic detectors}
\author{}\date{16 November 2026}
\begin{document}\maketitle
\begin{abstract}
We give revised general and conditional elliptic detectors while retaining
all point coordinates, comparison scales, and arithmetic height arguments.
In the general source, the physical weight becomes
$w=\rho R/(1+\rho R)^2$ and each scalar numerator gains one factor of $R$.
The scalar extraction argument then gives the required divisibility in
terms of $|v_p(\rho)|$. In the conditional source, the six factors are
$1+C_i^2/F$; an invertible change of the last comparison makes the product
symmetric under the reference involution. Ordered completeness and
whole-character guards hold for both revisions. Under the retained
arithmetic premises they give twenty-five and thirteen total variables,
respectively, including the input. No numerical determinant degree bound
or smaller defining-variable bound is asserted.
\end{abstract}

\section{The revised general source}

Fix the general twenty-five-variable source with unsquared comparison
scale in~\cite{Unit}. Its scalar and input label may be either those
of that source or the detector-field norm scalar of~\cite{Aligned}.
Write the integer-reflecting input label as $\beta(X)$, the scalar as
$\rho_L(X)$, and the consumed prime as $\ell=76249$. The exponent
$L\ge3$ is fixed and odd. All assertions about the integer input
projection retain the source's specified group, index-height,
normalized contact-height, pole-realization, prime-pattern and
aggregate trace-height premises~\cite{General,Unit,Aligned}.

Use the fixed curve and point
\[
 K=\Q(\sqrt2),\quad D=(3-\sqrt2)^2,\quad E:y^2=x^3-Dx,
 \quad P=\left(\frac{57-28\sqrt2}{8},
                    \frac{-660+417\sqrt2}{32}\right).
\]
The exceptional set $S$, nine point roles, comparison orders, and
all native point and difference-chart guards are retained. Put
\begin{equation}\label{eq:contacts}
\begin{gathered}
 a_0=1,\quad a_1=1+\rho T^2,\quad t_j=a_j\lambda,\quad
 \delta_j=t_j-\gamma_j,\quad\gamma_j=-j-1,\\
 \Theta=\{0,-1/2,-1,1/2399,-32/407\},\quad
 \Theta_j=(4j+2+\Theta)\cup(4j+4+\Theta),\\
 R_j=\frac{\prod_{\theta\in\Theta_j}(t_j-\theta)}{\delta_j^{10}},
 \qquad u_j=\rho R_j,\qquad
 w_j=\frac{u_j}{(1+u_j)^2}\qquad(j=0,1).
\end{gathered}
\end{equation}
Here $T$ is the retained rational trace of the sum of the four root
abscissas. Adjoin the scalar guards $1+u_j\ne0$.

For clarity, the nine point roles are
$A,S,T_1,Z,H,U_A,U_S,U_T,U_Z$. Set
$Q_1=A,Q_2=S,Q_3=T_1,Q_4=Z$, $U_0=H$ and
$(U_1,U_2,U_3,U_4)=(U_A,U_S,U_T,U_Z)$. With
$f=1/x$, $z=-x/y$, write
\[
 V_0=H-P,\quad V_i=U_i-Q_i,\quad
 L_i=z(U_i)/z(H),\quad
 L'_i=\overline{z(V_i)/z(V_0)}.
\]
The sixteen field tests, in order, are
\[
 (f(U_i))_{i=0}^{4},\quad
 (\overline{f(V_i)})_{i=0}^{4},\quad
 (L_i-L'_i)_{i=1}^{4},\quad
 L_1-\beta,\quad \mathcal F(L_2,L_3,L_4),
\]
where
$\mathcal F(a,t,b)=t^2-a(a-2)tb+(a-2)^2b^2-1$.
Their orders $a_i$ are one except at positions $12,13,14,16$,
which have order three. The six rational scalar tests $S_k$ are
the two components of each of
\[
                     f(H),\qquad f(S-P),\qquad z(S-P)/z(H),
\]
with orders $\kappa_k=(1,1,1,1,3,3)$.
The scalar normalizer acquires one additional factor of $R_j$:
\[
 \Phi_R(s)=\frac{Rs}{1+Rs},\qquad
 G_{j,k}=R_j^{\kappa_k+1}\Phi_{R_j}(S_k).
\]
The revised detector is
\begin{equation}\label{eq:detector}
\begin{gathered}
 W_j=\prod_{i=1}^{16}
   \left(1+\frac{C_i^2}{\delta_j^2 w_j^{2a_i-1}}\right)
   \prod_{k=1}^{6}
   \left(1+\frac{G_{j,k}^2}{w_j^{2\kappa_k-1}}\right),\\
 V_j=\Tr_{K/\Q}(W_j^{-1}),\qquad
 J_j=\Nm_{K/\Q}(W_j^{-1}),\qquad
 b_j=\frac{4w_j}{V_j^2},\qquad K=\Q(\sqrt2).
\end{gathered}
\end{equation}
Thus the new physical weight is used in every occurrence in
\eqref{eq:detector}. The scalar $\rho$, the contact functions $R_j$,
and the comparison scales remain those in~\eqref{eq:contacts}.
All displayed divisions retain their nonzero guards.

\begin{theorem}\label{thm:source}
Under the retained arithmetic premises, the source
\eqref{eq:contacts}--\eqref{eq:detector}, with the two native pole
blocks, has integer input projection for every sufficiently large
fixed odd $L$. Completeness holds for every integer input and every
fixed odd $L\ge3$.

Its finite-algebra determinant has twenty-five total variables.
Every odd valuation of each revised physical weight is positive.
With the scalar of~\cite{Aligned}, both prepared physical weights
can still be chosen as norms from $K$. No numerical sufficient value
of $L$ or determinant degree bound is asserted for this specification.
\end{theorem}

\section{Local extraction with absolute scalar valuations}

\begin{lemma}\label{lem:weight}
For $u\in\Q\setminus\{0,-1\}$ put $w=u/(1+u)^2$. If
$n=v_p(u)$ is odd, then
\[
                              v_p(w)=|n|>0.
\]
If $n$ is even, $v_p(w)$ is even. For $u>0$ one has
$0<w\le1/4$. Multiplication by the displayed rational square
preserves the square class and membership in every quadratic norm
subgroup of $\Q^\times$.
\end{lemma}
\begin{proof}
For $n>0$ the denominator is a unit; for $n<0$ it has valuation
$2n$. If $n=0$, the weight valuation is $-2v_p(1+u)$.
These cases prove the parity assertions. The real bound follows
from $(1+u)^2-4u=(1-u)^2\ge0$. A rational square is a norm
from every quadratic field.
\end{proof}

\begin{lemma}\label{lem:trace}
Let $p\notin S$ and suppose $\alpha=v_p(w_j)$ is odd. If $b_j$
has no negative odd valuation, then at some place above $p$ every
factor of $W_j$ has valuation greater than $-\alpha/2$. At that place,
\[
                 v(G_{j,k})>(\kappa_k-3/4)\alpha.
\]
If $\delta_j$ is a unit, one also has
\[
                         v(C_i)>(a_i-3/4)\alpha.
\]
\end{lemma}
\begin{proof}
By Lemma~\ref{lem:weight}, $\alpha>0$. Each nonzero second
summand of a detector factor has odd valuation, so each factor has
nonpositive valuation. Put
$q=\max_{\mathfrak p\mid p}v_{\mathfrak p}(W_j)\le0$.
Integrality of local traces gives $v_p(V_j)\ge-q$.
Since $v_p(b_j)$ is odd and not negative,
\[
                        1\le v_p(b_j)\le\alpha+2q.
\]
Thus $q>-\alpha/2$. At a place attaining $q$, each factor
valuation is at least their sum $q$. Applying this to its odd-order
second summand gives the stated inequalities. Zero tests contribute
the factor one and satisfy the inequalities automatically.
\end{proof}

Let $c=|n_S-1|$ be the nonzero adjacent root-index difference in the
source, with $S=[n_S]P$. The letter $S$ for this point is independent
of the exceptional prime set.

\begin{proposition}\label{prop:scalar}
Let $p\notin S$ and suppose $\mu=v_p(\rho)$ is odd. In a source
solution at least one $R_j$ is a unit, and
\begin{equation}\label{eq:scalarorders}
\begin{gathered}
 v_p(S_k)>|\mu|/4\quad(1\le k\le4),\qquad
 v_p(S_k)>9|\mu|/4\quad(k=5,6),\\
 v_p(c)>9|\mu|/4,\qquad
                     v_p(c)\ge2|\mu|+1\ge3.
\end{gathered}
\end{equation}
No positivity assumption on $\mu$ is needed for these conclusions.
\end{proposition}
\begin{proof}
The scale $a_1=1+\rho T^2$ and the two contact sets have not changed.
The original disjoint-contact proof therefore still gives a unit
$R_j$: $a_1$ is congruent to one or has negative valuation, and an
argument of positive or negative valuation gives a unit in a balanced
contact function. If the first contact function is not a unit at a
unit argument, the second contact-and-pole set is disjoint from its
residue~\cite[Section 3]{General}.

At this block, Lemma~\ref{lem:weight} gives
$\alpha=v_p(w_j)=|\mu|$. Apply Lemma~\ref{lem:trace}.
Since $R_j$ is a unit, positive valuation of $G_{j,k}$ is equivalent
to positive valuation of $S_k$, and the two valuations are equal.
This proves the first two groups of inequalities.

The first four scalar components put $H=[h]P$ and $S-P$ in the
formal group at every place above $p$. The unchanged formal logarithm
identity is
\[
 v_{\mathfrak p}\!\left(\frac{z(S-P)}{z(H)}\right)
                    =v_p((n_S-1)/h).
\]
The last two scalar components give the lower bound
$9|\mu|/4$ on the left. Since $h$ is a nonzero integer,
$v_p(h)\ge0$. This proves the bound for $c$. Integral valuations
and $|\mu|\ge1$ give the last assertions.
\end{proof}

\begin{proposition}\label{prop:mixed}
Let $p\notin S$ and suppose $\nu=v_p(R_j)>0$ is odd.
A source solution has either
\[
 v(C_i)\ge a_i\quad\hbox{for every $i$ at one place above $p$},
 \qquad\hbox{or}\qquad v_p(c)\ge3.
\]
\end{proposition}
\begin{proof}
The normalizer $\delta_j$ is a unit at such a contact. If
$\mu=v_p(\rho)$ is odd, Proposition~\ref{prop:scalar} gives the
second alternative. If $\mu$ is even, $\mu+\nu$ is odd and
\[
                 v_p(w_j)=|\mu+\nu|\ge1.
\]
The trace lemma gives
$v(C_i)>(a_i-3/4)|\mu+\nu|\ge a_i-3/4$.
Integral valuations give the first alternative.
\end{proof}

\section{Heights and soundness}

Only the conclusions of Propositions~\ref{prop:scalar}
and~\ref{prop:mixed} are used in the source's subsequent residual
and contact-height arguments. In particular, the original elimination
of negative odd scalar valuations is not needed separately.

For the original scalar, every finite odd contact of the label outside
the fixed exceptional set has positive odd scalar valuation. For the
scalar of~\cite{Aligned}, the same statement holds for its rational
contact argument after omitting the fixed numerator primes specified
there. Negative label-denominator valuations do not contribute odd
contacts outside that fixed set. Proposition~\ref{prop:scalar}
therefore gives the same bound
\[
                         h(\beta)\le Cc^{7/9}.
\]
All omitted primes and the fixed affine rescaling of the contact
argument affect only the fixed constant $C$.

The consumed-prime formula is unchanged because $\rho$ is unchanged.
It has positive odd valuation at least $L$, so
\[
                 c\ge\ell^{2L+1},\qquad
                 L\log\ell\le\tfrac12\log c.
\]
The scalar height estimate also remains the one from the chosen source:
\[
                  h(\rho)\le C_0h(\beta)+2L\log\ell+C_1,
\]
where $C_0,C_1$ are fixed. The comparison scale is still
$a_1=1+\rho T^2$, so its trace-height inequality is unchanged.

Proposition~\ref{prop:mixed} gives exactly the original contact
radical divisibilities for the input residual and the Lucas residual.
The balanced contact functions and their two height estimates are
unchanged. Consequently the unsquared-source proof in~\cite{Unit}
gives
\[
 B_{\mathrm{ind}}^2\le C_2|Nc|^\eta
             +C_2c^{7/9}+C_2\log c+C_3,
\]
with the same exponents $\eta=7/18$ and $\eta=7/6$ in the two
residual cases. Its strict exponent comparisons force the Lucas
residual to vanish and then force $\beta(X)$ to equal an integer
root label for every sufficiently large fixed odd $L$.
The selected label polynomial reflects integrality, so $X\in\Z$.
This proves soundness in Theorem~\ref{thm:source}.

\section{Ordered completeness and whole-character guards}

Fix an integer input and an odd $L\ge3$. The chosen scalar is a
positive prepared rational number. Its limiting physical value is
\[
                         w_\infty=\frac{\rho}{(1+\rho)^2}>0.
\]
It differs from $\rho$ by a rational square. Enlarge the finite
preparation and avoidance lists by the prime support of $1+\rho$.
This is permitted by the retained prime-pattern premise.

Choose the shared reference point, the adjacent root label, the Lucas
root labels, and the scalar neighborhoods in the original order.
The six scalar components can be made arbitrarily small at all the
fixed finite places and at the real places. With $R_j$ close to one,
$w_j$ is close to the fixed nonzero value $w_\infty$, and
$G_{j,k}=R_j^{\kappa_k+1}\Phi_{R_j}(S_k)$ tends to zero
as $S_k$ tends to zero. Thus every scalar
second summand in~\eqref{eq:detector} can be made as small as required.
No slope is selected at this stage.

Fix the root points and scales, avoid the finitely many point and
difference poles, and choose compact prepared real neighborhoods and
sufficiently small finite-place neighborhoods. All comparison
functions then have fixed bounds there. Apply the original simultaneous
prime-pattern theorem with its fixed avoidance, denominator divisibility,
large positive real slope, and conjugate reduction-order conditions.
As the slope tends to infinity at a preparation place,
\[
                     R_j\longrightarrow1,\qquad
                     w_j\longrightarrow w_\infty,
                     \qquad\delta_j^{-1}\longrightarrow0.
\]
Thus the comparison factors tend to one, uniformly on the already fixed
neighborhoods. All prescribed finite precisions and real bounds can be
imposed before prime selection. The common-index congruence refinement
and the additional selected-prime divisibility are unchanged.

At a selected fresh comparison prime, $\rho,\delta_j$ are units and
$v_p(R_j)=1$. Thus $v_p(w_j)=1$, and the original exact multiplication
orders give successful-slot comparison factors congruent to one.
The original local argument gives $v_p(\Phi_{R_j}(S_k))\ge0$.
The additional power of $R_j$ therefore gives
$v_p(G_{j,k})\ge\kappa_k+1$, so each scalar second summand
has valuation at least three at both places. The reciprocal product is
one modulo the selected place, while at the conjugate place it is
zero or one. The reciprocal trace is a unit, so the odd valuation
of $b_j$ is positive.

Every other odd valuation of $w_j$ is supported on an odd scalar prime.
For integer inputs those original scalar valuations are positive.
The finite preparation makes both products sufficiently close to one
there, preserving a unit trace. All remaining physical valuations are
even and remain even after division by the square trace. Thus the
reduced denominator of $b_j$ is a square.

At $2,3,17$, the scalar has positive valuation, so $1+\rho$ is a unit.
The limiting physical weight has the same local square class and
valuation as $\rho$. Taking the slope and comparisons sufficiently
close to their limits gives the required local squareclasses and
valuation depths for the prepared-input pole theorem. At real places,
choose $R_j>0$ and $1\le W_j<2$ in both slots. Then
\[
                    0<w_j\le\frac14,\qquad
                    V_j>1,\qquad 0<b_j<1<17.
\]
The retained pole-realization theorem supplies both pole blocks.

For every geometric character of the point algebra, the source's point
and difference guards are retained by the same nonzero-label and
sign-branch argument of~\cite[Section 5]{Unit}. The scalar normalizer
zero sets are avoided in the same prepared neighborhoods. The new
factors $1+u_j$ are positive rational scalars at these tuples, since
$\rho>0$ and $R_j>0$; hence they are whole-algebra units.
The additional power of $R_j$ introduces no denominator.
The revised detector factors remain one plus a nonnegative real square
times a positive coefficient at every character. All product and trace
norms are therefore units. The original density argument for the pole
parameters avoids every remaining conjugate pole denominator and
radicand.

No rational coordinate was added. The retained coordinates remain the
input, eighteen point coordinates, one slope and four pole coordinates;
the finite-algebra separator adds one coordinate. This proves the
count and completeness assertion in Theorem~\ref{thm:source}.
With the norm scalar, the prepared $u_j$ are norms from $K$ by
\cite{Aligned}. Lemma~\ref{lem:weight} preserves this property.

\section{The revised conditional source}

Retain the three point roles $A,H,U$, the common comparison slope,
and the scales
\[
 \beta=X^2+B,\quad T=\Tr(x(A)),\quad a_0=1,\quad
 a_1=\beta^d+T^d,\qquad (\kappa,d)\in
 \{(121/100,6),(16/15,4)\}.
\]
Here $B\ge3$ is a fixed odd integer. The weak rational $abc$ hypothesis
at exponent $\kappa$, together with the other arithmetic premises,
has exactly the scope specified in~\cite{Conditional,Unit}.
Put $f=1/x$, $z=-x/y$, and write
\[
 r=\frac{z(U)}{z(H)},\qquad
 s=\frac{z(U-A)}{z(H-P)}.
\]
Use the six tests
\begin{equation}\label{eq:conditionaltests}
 \begin{gathered}
 C_1=f(H),\quad C_2=f(U),\quad
 C_3=\overline{f(H-P)},\quad C_4=\overline{f(U-A)},\\
 C_5=r-\bar s,\qquad C_6=r+\bar s-2X.
 \end{gathered}
\end{equation}
Compared with the pair $r-\bar s,r-X$, the transformation on the last
two tests has matrix
\[
                         \begin{pmatrix}1&0\\-1&2\end{pmatrix}.
\]
Its determinant is a unit at every nonexceptional prime. For every
integer local threshold $m$, the old pair has both valuations at least
$m$ if and only if the new pair does. The ratio $\bar s$ already
occurs in the old fifth test, so no new chart denominator is introduced.

At each argument $t=a_j\lambda$ put
\begin{equation}\label{eq:conditionaldetector}
 \begin{gathered}
 F(t)=\frac{t(t+1)(2t+1)(3t+1)(4t+1)(5t+2)}{120},\qquad
 \delta=t-1,\qquad w=\frac{F(t)}{4\delta^8},\\
 W=\prod_{i=1}^6\left(1+\frac{C_i^2}{F(t)}\right),\qquad
 V=\Tr(W^{-1}),\qquad b=\frac{4w}{V^2}.
 \end{gathered}
\end{equation}
Retain the original field and scalar guards and the two pole blocks.
In particular $F$, $\delta$, $\Nm(W)$ and $V$ are nonzero.
The physical weight, its six contacts, and its pole of even order
are unchanged. The factor numerators no longer contain $\delta^4$.

\begin{lemma}\label{lem:conditionalodd}
For every guarded rational $t$, every odd valuation of $w(t)$
outside the conditional exceptional set is positive. At a positive
odd contact $\nu=v_p(w)$, absence of a negative odd valuation of $b$
forces all six tests in~\eqref{eq:conditionaltests} to have positive
valuation at one place above $p$.
\end{lemma}
\begin{proof}
All denominators, contact collisions, and the value $F(1)=7$ are units
outside the exceptional set. If $v_p(t)<0$, then
$v_p(F)=6v_p(t)$ and $v_p(\delta)=v_p(t)$, so
$v_p(w)=-2v_p(t)$ is even. If $v_p(\delta)>0$, then
$v_p(w)=-8v_p(\delta)$ is even. In the other cases $w$ is
integral, proving the first assertion.

At an odd contact, $\delta$ is a unit, $v_p(F)=\nu>0$, and each
factor in~\eqref{eq:conditionaldetector} has valuation
\[
                         \min\{0,2v(C_i)-\nu\}.
\]
The minimum is unique. The same trace argument as in
Lemma~\ref{lem:trace} gives a place at which every factor valuation
exceeds $-\nu/2$. Hence $v(C_i)>\nu/4>0$ for all six tests.
The inverse comparison transformation recovers the same conclusion
for the original target test $r-X$.
\end{proof}

\begin{theorem}\label{thm:conditionalsource}
Under the retained conditional arithmetic premises, for each displayed
pair $(\kappa,d)$, the revised source defines the integers over the
rationals for every sufficiently large fixed odd $B$. Completeness
holds for every nonzero integer input and every odd $B\ge3$; zero
is included by the source's final input factor. The finite-algebra
polynomial has thirteen total variables. No numerical sufficient
shift or determinant degree bound is asserted here.
\end{theorem}
\begin{proof}
For soundness, the local proof in~\cite[Section 4]{Conditional} uses
only positivity of the valuations of the six tests at the successful
place. Lemma~\ref{lem:conditionalodd} and the inverse comparison
transformation give those same conditions. The first four tests put
the two pairs of points in the indicated formal groups; the two ratio
tests recover the integer root label modulo every consumed contact
prime. Thus the reduced numerator of $X-n$ is divisible by the same
odd contact radicals. The two contact-height estimates, root-height
bound, comparison scale and sufficient-shift contradiction are unchanged.
They prove soundness for sufficiently large fixed odd $B$.

For completeness, fix $X=n\ne0$, take $A=[n]P$, and temporarily
put $U=[n]H$. Choose the source's compact real neighborhood and fixed
finite-place neighborhoods for $H$, with all point and difference
functions defined and bounded. The new comparison is a fixed linear
combination of the old ones, so it is bounded on the same neighborhoods.
The source permits any additional fixed precision at the exceptional
prime two.

Fix the roots, scales and neighborhoods before applying the original
prime-pattern theorem. It supplies the slope, the same six prime contacts,
and the compatible conjugate reduction orders. Refine the auxiliary
index while retaining the prescribed neighborhoods. At the successful
place of a selected contact prime the first four tests have valuation
at least two; the exact multiplication identities give valuation at
least four for both comparison tests. Thus each revised factor equals
one modulo $p^3$. At the opposite place every factor has nonpositive
valuation, so its reciprocal product reduces to zero or one. The trace
is a unit. All odd valuations of the physical input are positive.

At fixed finite preparation places, the slope can be made sufficiently
large in the local absolute value before the final congruence refinement.
The tests have fixed lower valuation bounds there, and
\[
 \frac1{F(t)}\sim t^{-6},\qquad
 w(t)\sim\frac1{4t^2},\qquad W=1+O(t^{-6}).
\]
Therefore the products can be made as close to one as required. The
original physical weight gives the same square reduced denominator and
the same prepared squareclasses and depths at $2,3,17$. At real places
these estimates hold uniformly on the chosen compact neighborhoods;
hence $0<b<17$ for sufficiently large positive slope. The selected
contact primes are the only possible new odd weight primes, exactly
as in the original prime-times-square identity. At ordinary poles
$t\equiv1$, the weight valuation is even, so no extra condition is
needed for the square reduced denominator.

The new factors have positive coefficient $1/F$ at every prepared real
character and are at least one. The all-character positivity and
point-chart argument of~\cite[Section 5]{Unit} therefore applies.
The comparison change adds no denominator, and the revised factor
coefficient uses only $F$, already guarded. The independent prepared
pole witnesses avoid every conjugate pole denominator and radicand
by the same density argument. The coordinate count is unchanged:
one input, six point coordinates, one slope, four pole coordinates,
and one finite-algebra separator.
\end{proof}

\begin{proposition}[Reference symmetry]\label{prop:symmetry}
On the common rational-function domain, the involution
\[
                         \tau(H,U)=(P-H,A-U)
\]
sends the product in~\eqref{eq:conditionaldetector} to its conjugate.
Consequently its reciprocal trace, reciprocal norm, and rational
functions of them are invariant under $\tau$.
\end{proposition}
\begin{proof}
The functions $f,z$ are respectively even and odd under point negation.
The involution interchanges $r$ and $s$. It sends the first four tests,
in order, to $\bar C_3,\bar C_4,\bar C_1,\bar C_2$, and sends the
last two to $-\bar C_5,\bar C_6$. Squaring removes the sign, proving
the product identity. The root point $A$ and hence the scales and $F$
are fixed by $\tau$.
\end{proof}
The identity is an equality of rational functions on the common domain;
it does not assert that a particular affine guarded chart is invariant
under the involution or that every rational point of its quotient lifts.

\section{Finite encoding and scope of the revisions}

Both revised constructions retain the same free point algebras, pole
parameters, finite quadratic pole roots, and one-coordinate separator
as their sources. Each added scalar denominator is guarded, and its
nonvanishing at every prepared character has been proved above.
Polynomial clearing, multiplication matrices, and norm guards therefore
produce integer polynomial specifications with the following counts:
\begin{center}
\begin{tabular}{lrrrrr}\toprule
Source&Input&Point coordinates&Slope&Pole coordinates&Separator\\\midrule
General&1&18&1&4&1\\
Either conditional&1&6&1&4&1\\\bottomrule
\end{tabular}
\end{center}
The totals are twenty-five and thirteen. All coefficients and
finite-algebra operations are computable after fixing the source
parameters. The sufficient defining parameters retain the source's
ineffectivity. Degree and coefficient bounds must be computed for the
new specifications; no previous numerical bound is transferred merely
from the unchanged coordinate count.

These revisions also give positive odd physical valuations and stronger
local preparation for further finite square conditions. A proof that
such conditions have simultaneous rational roots at the two actual
detector arguments is a separate arithmetic requirement. Neither the
source transfer, the reference symmetry, nor finite-place preparation
proves that requirement or a twenty-one- or nine-variable bound.

\begin{thebibliography}{9}
\bibitem{General}
\emph{A defining polynomial in twenty-five variables: combined scalar
and comparison conditions}, repository source,
\nolinkurl{papers/few_variables25/few_variables25_merged_20261110_v19.tex}.
\bibitem{Unit}
\emph{A single norm determinant for the 25- and 13-variable constructions},
repository source,
\nolinkurl{papers/finite_algebra/unit_norm_20261114_v28.tex}.
\bibitem{Aligned}
\emph{A quartic integral label and scalar norms from the detector field},
repository source,
\nolinkurl{papers/pole_parameters/detector_field_norm_scalar_20261116_v52.tex}.
\bibitem{Conditional}
\emph{Thirteen Variables from Direct Elliptic Index Comparison},
repository source,
\nolinkurl{papers/conditional_few_variables/direct_input13_20261110_v22.tex}.
\end{thebibliography}
\end{document}
